\bookStart{Högstena plate (Vg 216)}

\begin{flushright}%
\textbf{Dating:} C12th

\textbf{Meter:} \Fornyrdislag%
\end{flushright}%

\section{Introduction}

A runic-inscribed bronze plate found in a graveyard by Högstena church in Gudhem, West Geatland, Sweden.

NOTE: The original Norse text is a placeholder and should not be referenced.  Cf.~\textcite[169--188]{Pereswetoff2019}.

\sectionline

\section{Text}

\bvg\bva[]%
\alst{G}alanda viðr \hld\ \alst{g}ǫngla viðr &
\alst{r}íða\emph{n}da við\emph{r} \hld\ viðr \alst{r}innanda &
viðr \alst{s}ęljanda \hld\ viðr \alst{s}ign\emph{and}a &
viðr \alst{f}\emph{a}randa \hld\ viðr \alst{f}liuganda. &
Skal alt furna uk um døia.\eva

\bvb \evb\evg

\sectionline
